%% ============================================================
%%  PREAMBLE — Production-Grade Agentic AI: LangGraph + Python
%%  Extracted from main.tex.  Input ONCE from Tex/main.tex.
%% ============================================================

\documentclass[
  11pt,
  twoside,
  openright
]{book}

\usepackage{polyglossia}
\setdefaultlanguage{english}

\usepackage{fontspec}
\setmainfont{TeX Gyre Termes}[
  BoldFont   = TeX Gyre Termes Bold,
  ItalicFont = TeX Gyre Termes Italic
]
\setsansfont{TeX Gyre Heros}[
  BoldFont   = TeX Gyre Heros Bold,
  ItalicFont = TeX Gyre Heros Italic
]
\setmonofont{TeX Gyre Cursor}[
  Scale    = 0.88,
  BoldFont = TeX Gyre Cursor Bold
]

\usepackage{geometry}
\geometry{
  papersize  = {7in, 10in},
  top        = 1in,
  bottom     = 1.1in,
  inner      = 1in,
  outer      = 0.875in,
  headheight = 23.07pt,
  twoside
}

%% Fix fancyhdr \headheight warning (must be after geometry)
\setlength{\headheight}{23.07004pt}
\addtolength{\topmargin}{-0.00003pt}

\usepackage{xcolor}
\definecolor{BookBlue}       {HTML}{1A4A8A}
\definecolor{BookDark}       {HTML}{1C1C1E}
\definecolor{BookMid}        {HTML}{4A4A4A}
\definecolor{LearnBg}        {HTML}{EEF7EE}
\definecolor{LearnBorder}    {HTML}{27AE60}
\definecolor{TipBg}          {HTML}{EAF4FB}
\definecolor{TipBorder}      {HTML}{2980B9}
\definecolor{WarningBg}      {HTML}{FEF9E7}
\definecolor{WarningBorder}  {HTML}{F39C12}
\definecolor{NoteBg}         {HTML}{F4F4F4}
\definecolor{NoteBorder}     {HTML}{888888}
\definecolor{TakeawayBg}     {HTML}{F0F0FF}
\definecolor{TakeawayBorder} {HTML}{5B4FBE}
\definecolor{CodeBg}         {HTML}{F7F7F7}
\definecolor{CodeComment}    {HTML}{6A737D}
\definecolor{CodeKeyword}    {HTML}{D73A49}
\definecolor{CodeString}     {HTML}{032F62}
\definecolor{CodeConstant}   {HTML}{005CC5}
\definecolor{CodeAnnotation} {HTML}{6A737D}

\usepackage{microtype}
\usepackage{setspace}
\setstretch{1.18}

\usepackage{csquotes}
\usepackage{parskip}
\setlength{\parskip}{6pt}
\setlength{\parindent}{0pt}

\widowpenalty=10000
\clubpenalty=10000
\displaywidowpenalty=10000
\predisplaypenalty=10000
\raggedbottom
\emergencystretch=3em
\hbadness=3000
\vbadness=3000

\usepackage{needspace}
\usepackage{etoolbox}
\preto\section{\needspace{5\baselineskip}}
\preto\subsection{\needspace{4\baselineskip}}
\preto\subsubsection{\needspace{3\baselineskip}}

\usepackage{titlesec}
\renewcommand{\bottomtitlespace}{0.15\textheight}

\titleformat{\part}[display]
  {\centering\sffamily\Huge\bfseries\color{BookBlue}}
  {\sffamily\large\mdseries\partname\ \thepart}{20pt}{}

\titleformat{\chapter}[display]
  {\sffamily\huge\bfseries\color{BookBlue}}
  {\sffamily\normalsize\mdseries\chaptertitlename\ \thechapter}{12pt}{}
\titlespacing*{\chapter}{0pt}{0pt}{24pt}

\titleformat{\section}
  {\sffamily\Large\bfseries\color{BookBlue}}
  {\thesection}{1em}{}
\titlespacing*{\section}{0pt}{18pt}{6pt}

\titleformat{\subsection}
  {\sffamily\large\bfseries\color{BookDark}}
  {\thesubsection}{1em}{}
\titlespacing*{\subsection}{0pt}{14pt}{4pt}

\titleformat{\subsubsection}
  {\sffamily\normalsize\bfseries\color{BookDark}}
  {\thesubsubsection}{1em}{}
\titlespacing*{\subsubsection}{0pt}{10pt}{2pt}

\usepackage{fancyhdr}
\pagestyle{fancy}
\fancyhf{}
\fancyhead[LE]{\small\sffamily\color{BookMid}\leftmark}
\fancyhead[RO]{\small\sffamily\color{BookMid}\rightmark}
\fancyfoot[LE,RO]{\small\sffamily\color{BookMid}\thepage}
\renewcommand{\headrulewidth}{0.3pt}
\renewcommand{\footrulewidth}{0pt}

\fancypagestyle{plain}{
  \fancyhf{}
  \fancyfoot[LE,RO]{\small\sffamily\color{BookMid}\thepage}
  \renewcommand{\headrulewidth}{0pt}
}

\makeatletter
\renewcommand{\cleardoublepage}{%
  \clearpage
  \if@twoside
    \ifodd\c@page\else
      \thispagestyle{empty}\hbox{}\newpage
      \if@twocolumn\hbox{}\newpage\fi
    \fi
  \fi
}
\makeatother

\usepackage{rotating}
\usepackage{booktabs}
\usepackage{tabularx}
\usepackage{array}
\usepackage{longtable}
\renewcommand{\arraystretch}{1.35}

\newcolumntype{P}[1]{>{\RaggedRight\arraybackslash}p{#1}}

\usepackage{graphicx}
\usepackage{float}
\usepackage{caption}
\captionsetup{
  font      = {small,sf},
  labelfont = {bf,color=BookBlue},
  skip      = 6pt
}

\usepackage{listings}
\lstset{
  backgroundcolor   = \color{CodeBg},
  basicstyle        = \ttfamily\fontsize{9}{11}\selectfont,
  keywordstyle      = \color{CodeKeyword}\bfseries,
  stringstyle       = \color{CodeString}\itshape,
  commentstyle      = \color{CodeComment}\itshape,
  numbers           = left,
  numberstyle       = \tiny\color{BookMid},
  numbersep         = 6pt,
  breaklines        = true,
  frame             = single,
  framerule         = 0.4pt,
  rulecolor         = \color{NoteBorder},
  showstringspaces  = false,
  tabsize           = 2,
  captionpos        = b,
  breakatwhitespace = false,
  columns           = fullflexible,
  linewidth         = \linewidth,
  postbreak         = \mbox{\textcolor{gray}{$\hookrightarrow$}\space},
  xleftmargin       = 1.5em,
  framexleftmargin  = 1.5em,
  aboveskip         = 6pt,
  belowskip         = 4pt,
}

\lstdefinelanguage{Python}{
  morekeywords={
    False,None,True,and,as,assert,async,await,break,class,continue,def,
    del,elif,else,except,finally,for,from,global,if,import,in,is,lambda,
    nonlocal,not,or,pass,raise,return,try,while,with,yield
  },
  morekeywords=[2]{self,cls},
  sensitive=true,
  morecomment=[l]{\#},
  morestring=[b]",
  morestring=[b]',
  morestring=[s]{"""}{"""}, 
  morestring=[s]{'''}{'''},
  keywordstyle=\color{CodeKeyword}\bfseries,
  keywordstyle=[2]\color{CodeConstant},
  stringstyle=\color{CodeString},
  commentstyle=\color{CodeComment}\itshape,
  breaklines=true,
  breakatwhitespace=false,
  columns=fullflexible,
  keepspaces=true,
  showstringspaces=false,
  postbreak=\mbox{\textcolor{gray}{$\hookrightarrow$}\space},
}

\lstdefinelanguage{yaml}{
  morecomment=[l]{\#},
  morestring=[b]",
  morestring=[b]',
  literate=
    *{true}{{\color{CodeKeyword}true}}{4}
    {false}{{\color{CodeKeyword}false}}{5}
    {null}{{\color{CodeKeyword}null}}{4},
  sensitive=true,
  keywordstyle=\color{CodeKeyword}\bfseries,
  stringstyle=\color{CodeString},
  commentstyle=\color{CodeComment}\itshape,
  breaklines=true,
  columns=fullflexible,
  showstringspaces=false
}

\lstdefinelanguage{json}{
  morestring=[b]",
  morecomment=[l]//,
  morecomment=[s]{/*}{*/},
  sensitive=true,
  stringstyle=\color{CodeString},
  commentstyle=\color{CodeComment}\itshape,
}

\DeclareRobustCommand{\code}[1]{\texttt{#1}}
\providecommand{\term}[1]{\textit{#1}}
\providecommand{\chapterquote}[2]{%
  \vspace{1em}
  \begin{quote}\itshape#1\end{quote}
  \begin{flushright}--- #2\end{flushright}
  \vspace{1em}
}

\newcommand{\githubref}[1]{%
  \vspace{2pt}%
  {\small\sffamily\color{BookMid}Full source: \texttt{#1}}%
  \vspace{4pt}%
}

\usepackage{tcolorbox}
\tcbuselibrary{skins,breakable}

\newtcolorbox{learningbox}{
  colback=LearnBg, colframe=LearnBorder,
  before skip=12pt, fonttitle=\sffamily\bfseries\small,
  title={Learning Objective}, breakable, lines before break=3,
  left=8pt, right=8pt, top=6pt, bottom=6pt,
  boxrule=0.6pt, width=\linewidth
}
\newtcolorbox{tipbox}{
  colback=TipBg, colframe=TipBorder,
  before skip=12pt, fonttitle=\sffamily\bfseries\small,
  title={Tip}, breakable, lines before break=3,
  left=8pt, right=8pt, top=6pt, bottom=6pt,
  boxrule=0.6pt, width=\linewidth
}
\newtcolorbox{warningbox}{
  colback=WarningBg, colframe=WarningBorder,
  before skip=12pt, fonttitle=\sffamily\bfseries\small,
  title={Warning}, breakable, lines before break=3,
  left=8pt, right=8pt, top=6pt, bottom=6pt,
  boxrule=0.6pt, width=\linewidth
}
\newtcolorbox{notebox}{
  colback=NoteBg, colframe=NoteBorder,
  before skip=12pt, fonttitle=\sffamily\bfseries\small,
  title={Note}, breakable, lines before break=3,
  left=8pt, right=8pt, top=6pt, bottom=6pt,
  boxrule=0.6pt, width=\linewidth
}
\newtcolorbox{takeawaybox}[1][Key Takeaway]{
  colback=TakeawayBg, colframe=TakeawayBorder,
  before skip=12pt, fonttitle=\sffamily\bfseries\small,
  title={#1}, breakable, lines before break=3,
  left=8pt, right=8pt, top=6pt, bottom=6pt,
  boxrule=0.6pt, width=\linewidth
}

%% ── Source-Layout tree environment ─────────────────────────────────────
%% Usage:
%%   \begin{sourcelayout}
%%     \dirtree{%
%%       .1 src/.
%%       .2 agent/.
%%       .3 nodes/.
%%       .4 perceive\_node.py \DTcomment{§1.2}.
%%     }
%%   \end{sourcelayout}
%%
%% Requires: \usepackage{dirtree}  (loaded below)
%% Colours reuse NoteBg / NoteBorder (neutral grey — reference material,
%% not instructional content).
%% ────────────────────────────────────────────────────────────────────────
\usepackage{dirtree}

\newtcolorbox{sourcelayout}{%
  colback      = NoteBg,
  colframe     = NoteBorder,
  colbacktitle = NoteBorder,
  coltitle     = white,
  fonttitle    = \sffamily\bfseries\small,
  title        = {Source Layout},
  before skip  = 12pt,
  after skip   = 10pt,
  breakable,
  lines before break = 3,
  left   = 6pt,
  right  = 6pt,
  top    = 6pt,
  bottom = 6pt,
  boxrule = 0.6pt,
  width   = \linewidth,
}

%% dirtree global style: monospaced font, consistent with code listings
\renewcommand*\DTstylecomment{\rmfamily\color{CodeComment}\itshape}
\renewcommand*\DTstyle{\ttfamily\color{BookDark}}

\usepackage{amsmath}
\usepackage{amssymb}

\PassOptionsToPackage{hyphens}{url}
\usepackage{hyperref}
\usepackage{imakeidx}
\makeindex[columns=2, title=Index, intoc=true]
\urlstyle{tt}
\makeatletter
\g@addto@macro\UrlBreaks{\do\.\do\_\do\-}
\makeatother

\hypersetup{
  colorlinks  = true,
  linkcolor   = BookBlue,
  urlcolor    = BookBlue,
  citecolor   = BookBlue,
  pdftitle    = {Production-Grade Agentic AI: LangGraph + Python Frameworks},
  pdfauthor   = {Haythem Aber},
  pdfsubject  = {LangGraph, Python, Agentic AI, Production Systems},
  pdfkeywords = {LangGraph, Python, AI agents, production, observability, security}
}

\usepackage{tikz}
\usetikzlibrary{arrows.meta, positioning, shapes.geometric, shapes.misc, calc}
\usepackage{pdflscape}
\usepackage{ragged2e}

\newcommand{\component}[1]{\textsf{\textbf{#1}}}
\newcommand{\seechapter}[1]{\textit{(See Chapter~\ref{#1})}}
\newcommand{\seepart}[1]{\textit{(See Part~\ref{#1})}}

\newcommand{\partclose}[1]{%
  \cleardoublepage
  \vspace*{\fill}
  \begin{center}\sffamily\Large\color{BookBlue}#1\end{center}
  \vspace*{\fill}
  \cleardoublepage
}

%% Shared book-wide infrastructure layered on top of the author's preamble.
%% Keep this file loaded once by Tex/main.tex; chapter fragments must not load it.

\usepackage{calc} % Pandoc-generated longtable column dimensions.
\usepackage[
  backend=biber,
  style=verbose-note,
  autocite=footnote,
  sorting=nyt,
  maxbibnames=99
]{biblatex}
\addbibresource{Tex/bibliography/references.bib}

% Paper labels remain explicit and are not interchangeable.
\newcommand{\CoALACite}[1]{\footnote{\textsc{[CoALA]} See Sumers et al., \emph{Cognitive Architectures for Language Agents}, #1.}}
\newcommand{\FanLanCite}[1]{\footnote{\textsc{[Fan \& Lan]} See Fan and Lan, \emph{From Cognitive Architectures to Language Agents}, #1.}}

% The source paper is the theoretical basis in Parts I and II; examples and
% engineering decisions are explicitly identified in the prose.
\newcommand{\BookExampleLabel}{\textsc{[Book example]}}
\newcommand{\EngineeringLabel}{\textsc{[Engineering decision]}}

\usepackage[toc,nonumberlist]{glossaries}
\makeglossaries
\renewcommand*{\glossarypreamble}{\begin{multicols}{2}\raggedcolumns\small}
\renewcommand*{\glossarypostamble}{\end{multicols}}
\newglossaryentry{languageagent}{name={language agent},description={An agentic system that uses a language model as a core computational component and interacts with an environment}}
\newglossaryentry{cognitivelanguageagent}{name={cognitive language agent},description={A language agent that also uses internal processes to manage its information and behavior}}
\newglossaryentry{workingmemory}{name={working memory},description={Active information maintained for the current decision cycle}}
\newglossaryentry{longtermmemory}{name={long-term memory},description={Memory modules that retain information beyond the current decision cycle}}
\newglossaryentry{episodicmemory}{name={episodic memory},description={Memory of experience from earlier decision cycles}}
\newglossaryentry{semanticmemory}{name={semantic memory},description={Memory of knowledge about the world and the agent itself}}
\newglossaryentry{proceduralmemory}{name={procedural memory},description={Knowledge about how to perform actions and decision procedures, represented implicitly in model weights or explicitly in agent code}}
\newglossaryentry{grounding}{name={grounding},description={Procedures that execute external actions and process environmental feedback}}
\newglossaryentry{decisionprocedure}{name={decision procedure},description={The top-level procedure that coordinates the agent's repeated decision cycles}}


% Reusable review and recap blocks used throughout the book.
\newtcolorbox{chapterrecap}{
  colback=TakeawayBg, colframe=TakeawayBorder,
  before skip=10.5pt, fonttitle=\sffamily\bfseries\small,
  title={Chapter recap}, breakable, lines before break=3,
  left=8pt, right=8pt, top=5.25pt, bottom=5.25pt,
  boxrule=0.6pt, width=\linewidth
}
\newtcolorbox{reviewquestions}[1]{
  colback=NoteBg, colframe=NoteBorder,
  before skip=10.5pt, fonttitle=\sffamily\bfseries\small,
  title={#1}, breakable, lines before break=3,
  left=8pt, right=8pt, top=5.25pt, bottom=5.25pt,
  boxrule=0.6pt, width=\linewidth
}
\providecommand{\tightlist}{}

% Stable source references for implementation listings. The QR payload, the
% clickable source ID and "Open exact source" all consume the same URL argument.
% The visible source path/member and linked test path remain useful in print.
\newcommand{\BookSourceReferenceText}[9]{%
  {\sffamily\fontsize{8}{9}\selectfont\raggedright
    \setlength{\parskip}{0pt}%
    \href{#8}{\textbf{Source #1}}\hspace{0.55em}\textbar\hspace{0.55em}\texttt{#2}\par
    \textbf{Path:} \href{#8}{\path{#3}}\par
    \textbf{Member:} #4\hspace{0.5em}\textbf{Source lines:} #5\par
    \textbf{Test file:} \href{#9}{\path{#6}}\par
    \textbf{Tests:}\par \href{#9}{#7}\par
    \href{#8}{\textbf{Open exact source}}%
  }%
}

\newcommand{\BookSourceReferenceNoQR}[9]{%
  \par\smallskip\noindent
  \begin{minipage}{\linewidth}%
    \BookSourceReferenceText{#1}{#2}{#3}{#4}{#5}{#6}{#7}{#8}{#9}%
  \end{minipage}%
  \par\smallskip
}

\newcommand{\BookSourceReferenceQRContent}[9]{%
  \par\smallskip\noindent
  \begin{minipage}[c]{0.785\linewidth}%
    \BookSourceReferenceText{#1}{#2}{#3}{#4}{#5}{#6}{#7}{#8}{#9}%
  \end{minipage}%
  \hfill
  \begin{minipage}[c]{0.18\linewidth}%
    \centering
    % The source QR is a vector PDF asset generated with qrcode.sty.
    % Embedding it as a graphic lets hyperref create a real clickable area
    % around the full printed symbol (including its four-module quiet zone).
    \href{#8}{\mbox{{\setlength{\fboxsep}{0pt}\colorbox{white}{\includegraphics[height=\BookSourceQRHeight]{Tex/diagrams/source-qr/#1.pdf}}}}}\par
    {\scriptsize\sffamily Scan to inspect}%
  \end{minipage}%
  \par\smallskip
}

% Standard source references use a 19.5 mm QR including the quiet zone.
% Very dense commit-pinned URLs may opt into the 21 mm variant.
\newcommand{\BookSourceReferenceQR}[9]{%
  {\def\BookSourceQRHeight{19.5mm}\BookSourceReferenceQRContent{#1}{#2}{#3}{#4}{#5}{#6}{#7}{#8}{#9}}%
}
\newcommand{\BookSourceReferenceQRLong}[9]{%
  {\def\BookSourceQRHeight{21mm}\BookSourceReferenceQRContent{#1}{#2}{#3}{#4}{#5}{#6}{#7}{#8}{#9}}%
}


% Book-level metadata overrides the sample metadata retained in the supplied preamble.
\hypersetup{
  pdftitle={Reliable AI Agents: Architecture and Implementation},
  pdfauthor={Haythem Aber},
  pdfsubject={CoALA, cognitive architectures, agentic systems, Java production engineering},
  pdfkeywords={CoALA, Fan and Lan, cognitive architectures, agentic AI, Java, supply-chain disruption management, production systems}
}

\hypersetup{
  pdftitle={Chapter 13 Outline Review: Understanding Typed Impasses, Substates and Compilation: B2},
  pdfsubject={Proposed detailed outline for individual author approval},
  pdfkeywords={Fan and Lan, Soar, typed impasses, substates, chunking, B2, recovery}
}
\begin{document}
\thispagestyle{empty}
\begin{center}
{\sffamily\Large\color{BookBlue} FROM COGNITIVE ARCHITECTURES TO PRODUCTION AGENTIC SYSTEMS\par}
\vspace{1.2cm}
{\sffamily\huge\bfseries\color{BookBlue} Chapter 13 Outline\par}
\vspace{0.45cm}
{\sffamily\LARGE Understanding Typed Impasses, Substates and Compilation: B2\par}
\vspace{0.8cm}
{\large Part II -- From Historical Cognitive Architectures to Mechanism-Level Understanding\par}
\vfill
\begin{notebox}
\textbf{Review status:} Proposed outline v1.0; awaiting the author's explicit approval. No Chapter 13 prose has been drafted or authorized.

\textbf{Primary theoretical source:} Fan and Lan (2026), \emph{From Cognitive Architectures to Language Agents: A Mechanism-Level Review of Lineage, Convergence, and Migration Gaps}, arXiv:2607.23942v1. The chapter focuses on the Soar mechanism in Section III-A, the modern mapping in Section V-A--V-B, and the B2 residual hypothesis and proposed evaluation in Section VI and Appendix D.

\textbf{Historical cross-check:} The Soar Architecture documentation cited as reference [9] by Fan and Lan. It is used to verify native Soar semantics, not as a separate theoretical framework for the book.

\textbf{Approval rule:} Chapter prose begins only after this outline is individually approved. Approval of the Part II batch schedule does not approve this outline.
\end{notebox}
\vfill
{\sffamily\large Outline review version 1.0\par}
{\small 14 November 2026\par}
\end{center}
\clearpage
\setcounter{secnumdepth}{0}
\tableofcontents
\cleardoublepage
\subsection{1. Chapter purpose and role in the book}\label{chapter-purpose-and-role-in-the-book}

Chapter 12 examined how memory accessibility, retrieval latency, downstream utility and cost might be coupled at context assembly. Chapter 13 turns to a different control problem: what should happen when an agent cannot continue a decision or execution path safely, and how can the system recover without corrupting the state that made the original task meaningful?

The chapter is a standalone deep-dive into \textbf{B2}, the typed impasse--substate--compilation bundle. Its historical anchor is Soar. Soar does not treat every difficulty as an undifferentiated error message. Its operator-selection process can expose a typed impasse; the architecture creates a child substate whose purpose is to resolve that impasse; the substate produces a result that allows the parent state to proceed; and, when the impasse is successfully resolved, chunking can compile the dependency conditions and result into a reusable production.

Fan \& Lan use that historical chain to formulate a bounded agent-runtime hypothesis. An execution ledger should classify a failure, preserve the relevant parent state, open a bounded recovery subgraph, return a typed resolution, and compile a reusable recovery rule only when the conditions that made the resolution valid are known. The chapter must distinguish the historical Soar mechanism from this proposed modern translation. In particular, a Soar superstate is not simply an immutable snapshot copied into a child process, and a Soar impasse is not synonymous with every failed tool call or exception.

The objective is to teach the reader why \textbf{failure representation, recovery scope, parent-state integrity, result return, and learning from a resolved recovery} must be understood as one control chain. A retry loop may recover once; a reflection note may influence a later prompt; a checkpoint may preserve state; and a repair policy may be selected by an error classifier. None of those pieces alone establishes the complete B2 coupling.

\subsubsection{Chapter thesis}\label{chapter-thesis}

\begin{quote}
A safe recovery mechanism must preserve the authority and relevant context of the parent task while it investigates a typed failure in a scoped child process. It must return a result whose type and applicability can be checked, and it should compile reusable knowledge only from the dependencies that made the resolution valid.
\end{quote}

\subsubsection{Learning outcomes}\label{learning-outcomes}

By the end of the chapter, a technically curious reader should be able to:

\begin{enumerate}
\def\labelenumi{\arabic{enumi}.}
\tightlist
\item
  Explain how Soar represents operator-selection impasses and why an impasse is a typed control condition rather than a generic exception.
\item
  Distinguish tie, conflict, constraint-failure, and no-change impasses, including the state-level and operator-level no-change cases.
\item
  Explain the relationship between a parent state (superstate), a child substate, subgoal processing, results, resolution, elimination, regeneration and dependency tracking.
\item
  Describe how Soar chunking uses the dependency conditions and result of a successfully resolved impasse to construct a production that can match a later similar state.
\item
  Distinguish Soar's native substate semantics from Fan \& Lan's proposed protected recovery subgraph for modern agent runtimes.
\item
  Compare the partial B2 baselines---PALADIN, AgentDebugX, Shepherd, AgentHER and Reflexion---without implying that any one of them implements the complete residual bundle.
\item
  Explain why a checkpoint, retry, recovery action, reversible branch, reflection record or learned feedback item is not by itself an implementation of B2.
\item
  Design a language-neutral recovery contract with a versioned parent state, bounded child recovery scope, typed resolution, explicit return validation and dependency-aware compilation.
\item
  Identify the safety risks of misdiagnosis, parent-state corruption, stale dependencies, irreversible side effects, overgeneralized compilation and repeated failure.
\item
  Reconstruct the B2 experimental protocol and its falsifiers, including state traces, outcome metrics, safety/efficiency measures and the required ablations.
\item
  State the boundaries between B2 and B1, B3, B4, B5 and B6 without teaching those later mechanism studies prematurely.
\end{enumerate}

\subsection{2. Organizing distinctions}\label{organizing-distinctions}

The chapter must repeatedly preserve the following distinctions. They are not stylistic details; they prevent the reader from equating mechanisms that have different triggers or control authority.

\begin{longtable}[]{@{}
  >{\raggedright\arraybackslash}p{(\columnwidth - 2\tabcolsep) * \real{0.5000}}
  >{\raggedright\arraybackslash}p{(\columnwidth - 2\tabcolsep) * \real{0.5000}}@{}}
\toprule\noalign{}
\begin{minipage}[b]{\linewidth}\raggedright
Distinction
\end{minipage} & \begin{minipage}[b]{\linewidth}\raggedright
Correct treatment
\end{minipage} \\
\midrule\noalign{}
\endhead
\bottomrule\noalign{}
\endlastfoot
\textbf{Soar impasse vs.~generic execution error} & A Soar impasse arises in operator selection when the current decision cannot proceed under the preference structure. A modern tool failure may be classified as a recovery trigger in the proposed B2 execution ledger, but it is not automatically a Soar impasse. \\
\textbf{Impasse detection vs.~recovery content} & The typed impasse identifies why the parent decision is blocked. Recovery operators or actions are generated and evaluated inside the child substate; they are not the impasse itself. \\
\textbf{Superstate vs.~immutable snapshot} & Soar represents the parent as the superstate and tracks dependencies between child results and superstate working-memory elements. Fan \& Lan's proposed runtime says to freeze the relevant parent state. A versioned, immutable snapshot is an engineering translation, not a literal claim about Soar's implementation. \\
\textbf{Substate vs.~ordinary checkpoint} & A checkpoint may preserve data for resume. A Soar-like substate additionally has a typed reason for creation, a parent/child relation, a local decision process and controlled result return. \\
\textbf{Resolved vs.~eliminated vs.~regenerated impasse} & A resolved impasse produces results that let the parent proceed and creates an opportunity for chunking. An impasse may instead become irrelevant because the situation changes, or be regenerated when dependencies change. Those outcomes do not license the same learning conclusion. \\
\textbf{Result vs.~arbitrary child output} & A Soar result is a working-memory element or preference linked to a superstate and survives removal of the child state. A proposed agent result envelope should make its type, evidence, parent version and applicability conditions explicit. \\
\textbf{Chunking vs.~reflection or replay} & Soar chunking compiles dependencies that were necessary for a resolved subgoal into a generalized production. Storing a failure summary, replaying a trajectory or writing a reflection may influence later behavior but does not automatically reproduce this dependency-based compilation. \\
\textbf{In-episode recovery vs.~general skill governance} & B2 is an in-episode recovery transition whose outcome may yield a reusable rule. B5 concerns general skill invocation/governance across tasks and is handled as the closed candidate in Chapter 17. \\
\textbf{Runtime insertion point vs.~complete implementation} & A state graph, checkpoint, interrupt, fork or replay boundary can host a B2 design. Its presence alone does not demonstrate typed diagnosis, protected recovery, typed return and safe compilation. \\
\end{longtable}

\subsection{3. Proposed chapter structure}\label{proposed-chapter-structure}

\subsubsection{Opening --- Chapter at a glance}\label{opening-chapter-at-a-glance}

\begin{itemize}
\tightlist
\item
  Start with a familiar but dangerous situation: an agent encounters a failed operation and can either retry blindly, abandon the task, or stop the parent task while it diagnoses the failure.
\item
  Use the Product X payment-service incident as the running \textbf{{[}Book example{]}}. The known starting fact is simply that a payment retry failed. Further details---such as a timeout with an uncertain settlement outcome---must be visibly identified as illustrative assumptions, not paper evidence.
\item
  State the central risk: a recovery attempt can be locally plausible and still violate the parent's goal, repeat an irreversible action, overwrite important state, or create a reusable rule that is unsafe in a different situation.
\item
  Preview the historical Soar chain, the modern partial baselines, and the B2 residual.
\item
  Present learning outcomes in the shared \texttt{learningbox}.
\end{itemize}

\subsubsection{13.1 The problem: failure is not yet a recovery policy}\label{the-problem-failure-is-not-yet-a-recovery-policy}

\paragraph{13.1.1 A failed operation does not tell the agent what to do next}\label{a-failed-operation-does-not-tell-the-agent-what-to-do-next}

\begin{itemize}
\tightlist
\item
  Explain why a raw exception such as \texttt{timeout}, \texttt{not\ found}, \texttt{permission\ denied} or \texttt{invalid\ state} is only evidence about an event, not a complete recovery decision.
\item
  Show how the same outward failure can have different causes and consequences. A timeout after a payment request, for example, may leave the settlement state unknown; automatically repeating the request could be unsafe.
\item
  Distinguish detection, attribution, selection of recovery, execution of recovery, validation of outcome and learning from the episode.
\end{itemize}

\paragraph{13.1.2 The parent task must remain meaningful while recovery is attempted}\label{the-parent-task-must-remain-meaningful-while-recovery-is-attempted}

\begin{itemize}
\tightlist
\item
  State the parent goal and its invariants before opening a recovery branch.
\item
  Explain the difference between preserving a record of what happened and preserving the state/authority needed to resume coherently.
\item
  Introduce the parent/child recovery analogy as an intuition only; defer exact terms until the Soar mechanism is taught.
\end{itemize}

\paragraph{13.1.3 From retry loops to a control chain}\label{from-retry-loops-to-a-control-chain}

\begin{itemize}
\tightlist
\item
  Compare a blind retry, diagnosis followed by retry, reversible branch, reflection memory and scoped substate at a high level.
\item
  Preview the five operations that B2 connects: classify, isolate, recover, return and compile.
\item
  Explain that the goal is not maximal automation. Some failure types should lead to escalation, abstention or human review rather than another tool call.
\end{itemize}

\subsubsection{13.2 Soar: the historical mechanism behind B2}\label{soar-the-historical-mechanism-behind-b2}

This is the chapter's central historical teaching section. It should be developed slowly from the operator decision process rather than beginning with a list of terminology. The Soar-specific descriptions below are checked against the architecture manual cited by Fan \& Lan.\autocite{soar2025architecture}

\paragraph{13.2.1 Working memory, operators and preferences}\label{working-memory-operators-and-preferences}

\begin{itemize}
\tightlist
\item
  Introduce Soar's current problem representation in working memory.
\item
  Explain operators as candidate state transitions: an operator proposes how the current state might change.
\item
  Explain preferences as the information used to select among proposed operators.
\item
  Show the decision point where one operator may be selected, or the preference structure may fail to determine a valid next step.
\item
  Do not re-teach the entire Soar architecture; explain only the state and control elements needed to understand impasses and substates.
\end{itemize}

\paragraph{13.2.2 What makes an impasse typed?}\label{what-makes-an-impasse-typed}

\begin{itemize}
\tightlist
\item
  Define the impasse as a state in the decision process, not just a descriptive label attached to a failed action.
\item
  Explain the Soar categories reported in the architecture documentation cited by Fan \& Lan:

  \begin{itemize}
  \tightlist
  \item
    \textbf{Tie impasse:} multiple acceptable operators remain tied because the preferences do not distinguish them.
  \item
    \textbf{Conflict impasse:} conflicting better/worse preferences prevent a consistent choice.
  \item
    \textbf{Constraint-failure impasse:} required/prohibited constraints conflict, leaving no legal selection under the current preference structure.
  \item
    \textbf{No-change impasse:} no new operator is selected during the decision procedure. Distinguish state no-change (no acceptable or required preference suggests an operator, or all acceptable candidates have been rejected) from operator no-change (an operator is selected but no additional productions match during application, or no operator is selected in the next decision phase).
  \end{itemize}
\item
  Explain that the type constrains what kind of problem the child process must solve. Do not suggest that all types are repaired by generating another preference; constraint failure and no-change conditions have different resolution paths.
\item
  Keep a short source note: Fan \& Lan's own summary emphasizes tie, conflict, rejection and missing choice; the fuller type distinctions are drawn from the Soar Architecture source it cites.
\end{itemize}

\paragraph{13.2.3 The substate: a child problem with a parent relation}\label{the-substate-a-child-problem-with-a-parent-relation}

\begin{itemize}
\tightlist
\item
  Explain that an impasse creates a new state/subgoal whose purpose is to resolve the impasse.
\item
  Introduce \textbf{substate} and \textbf{superstate} only after the reader understands the parent problem and the need for a child problem.
\item
  Explain that the initial child state includes a description of the impasse cause and a reference to the superstate.
\item
  Explain that the child can select and apply its own operators; additional impasses can create nested subgoals.
\item
  Explain that the top-level and nested states form a subgoal stack. Problem solving can remain active at multiple levels, even though attention tends to focus on the newest subgoal.
\item
  Clarify that the child is not an independent agent with an unrelated goal: its purpose is scoped by the parent impasse.
\end{itemize}

\paragraph{13.2.4 Results: what crosses back to the parent?}\label{results-what-crosses-back-to-the-parent}

\begin{itemize}
\tightlist
\item
  Define a subgoal result as a working-memory element or preference created in the substate and linked directly or indirectly to a superstate.
\item
  Explain why such results are not simply any text or arbitrary value returned by a child process.
\item
  Explain the role of support and persistence at a conceptual level: some results remain because they are connected to the higher-level problem; unrelated child-local structures disappear when the substate is removed.
\item
  Explain that an impasse is resolved when child processing creates results that allow the parent state to select a new operator.
\item
  Show the return path: the child contributes a result; the parent decision can proceed; the child state and its local non-result contents are removed.
\end{itemize}

\paragraph{13.2.5 Resolution is not the only way an impasse can disappear}\label{resolution-is-not-the-only-way-an-impasse-can-disappear}

\begin{itemize}
\tightlist
\item
  Distinguish \textbf{resolved}, \textbf{eliminated} and \textbf{regenerated} impasses.
\item
  Resolved: a child result enables progress in the parent; this is the learning-relevant path.
\item
  Eliminated: a higher-level change or new external input makes the old subgoal irrelevant; no chunking opportunity should be claimed from this event.
\item
  Regenerated: a dependency used by the child changes in the superstate; the old substate is removed and a new one is created so that child reasoning is consistent with the changed parent context.
\item
  Explain the Goal Dependency Set (GDS) as the source mechanism that records relevant superstate working-memory dependencies for substate results and supports regeneration when those dependencies change.
\item
  Emphasize that dependency tracking protects semantic validity; it is not simply a general-purpose version number or transaction log.
\end{itemize}

\paragraph{13.2.6 Chunking: compiling the conditions that made the resolution possible}\label{chunking-compiling-the-conditions-that-made-the-resolution-possible}

\begin{itemize}
\tightlist
\item
  Explain the learning opportunity after an operator impasse is successfully resolved.
\item
  Describe how Soar's chunking mechanism summarizes processing in the substate.
\item
  Explain the two sides of the learned production:

  \begin{itemize}
  \tightlist
  \item
    conditions: the relevant state elements from the superstate that, through the substate's production chain, enabled resolution;
  \item
    action: the working-memory element or preference that produced the successful result.
  \end{itemize}
\item
  Explain variablization/generalization: conditions and action are expressed so that the new chunk can match a similar future situation, rather than memorizing only the original concrete state.
\item
  Explain why this can prevent the same impasse from requiring the same deliberate subgoal work again.
\item
  Distinguish chunking from caching the entire trace, storing a natural-language lesson, or memorizing a successful action without its applicability conditions.
\item
  Discuss overgeneralization as a risk that motivates careful dependency conditions; keep the focus on the source mechanism, not a complete theory of Soar learning.
\end{itemize}

\paragraph{13.2.7 Worked Soar cycle}\label{worked-soar-cycle}

\begin{itemize}
\tightlist
\item
  Walk through a small decision: the parent state has two candidate operators whose preferences do not distinguish them; a tie impasse is created; the child substate evaluates a relevant condition; it produces a result linked to the parent; the parent can now select an operator; chunking can compile the relevant parent conditions and result action.
\item
  Use a second short counterexample where the parent context changes during subgoal processing and the child must be regenerated rather than returning a stale result.
\item
  Keep the illustration simple and explicitly label it a \textbf{{[}Book example{]}} unless the exact domain example is directly present in a cited Soar source.
\end{itemize}

\subsubsection{13.3 What changes when the mechanism moves into a language-agent runtime?}\label{what-changes-when-the-mechanism-moves-into-a-language-agent-runtime}

\paragraph{13.3.1 From operator-selection impasse to execution failure classification}\label{from-operator-selection-impasse-to-execution-failure-classification}

\begin{itemize}
\tightlist
\item
  Explain Fan \& Lan's translation: an execution ledger classifies a failure or attributed root cause and uses that type to choose a recovery transition.
\item
  State explicitly that a modern tool error is not literally a Soar impasse. The commonality is typed failure-specific control, not identical trigger semantics.
\item
  Explain why an untyped \texttt{failed} flag cannot choose safe recovery scope or compilation conditions.
\end{itemize}

\paragraph{13.3.2 From superstate/substate to protected parent and bounded recovery subgraph}\label{from-superstatesubstate-to-protected-parent-and-bounded-recovery-subgraph}

\begin{itemize}
\tightlist
\item
  Translate the parent/child relation into a proposed runtime structure: a versioned parent-state reference or snapshot, a child recovery scope, and a boundary around which mutations the child may make.
\item
  Mark the difference: Fan \& Lan's proposal says to freeze the relevant parent state; Soar's native semantics use a superstate relation and dependency tracking rather than a literal immutable cloned snapshot.
\item
  Explain why a generic fork or checkpoint is only a substrate. A B2-like branch needs a reason for its creation, a defined scope, an allowed mutation set and a return condition.
\end{itemize}

\paragraph{13.3.3 From Soar results to typed resolution contracts}\label{from-soar-results-to-typed-resolution-contracts}

\begin{itemize}
\tightlist
\item
  Translate the returned result into an explicit resolution object with a result type, evidence/provenance, parent version, relevant dependency set, postconditions and applicability boundary.
\item
  Explain why the parent should validate the resolution before resuming the original task.
\item
  Present the contract fields as \textbf{{[}Engineering decision{]}}, not as a schema prescribed by Fan \& Lan.
\end{itemize}

\paragraph{13.3.4 From chunking to condition-safe recovery-rule compilation}\label{from-chunking-to-condition-safe-recovery-rule-compilation}

\begin{itemize}
\tightlist
\item
  Explain the proposed compilation target: a narrow reusable recovery rule whose guard records the failure type and dependency conditions that made the original resolution valid.
\item
  Explain why a rule must not be generalized from an accidental correlation or a failed trajectory with ambiguous attribution.
\item
  Explain compile gating: only a successfully resolved and validated recovery can become a compilation candidate; unresolved, eliminated, regenerated or externally invalidated cases do not justify the same inference.
\item
  Explain versioning, provenance and invalidation as \textbf{{[}Engineering decision{]}} needed to keep a rule safe as tools, APIs, policies and environments change.
\end{itemize}

\subsubsection{13.4 Modern systems: the closest pieces are distributed}\label{modern-systems-the-closest-pieces-are-distributed}

The purpose of this section is not to write mini-surveys of every named system. It is to show which edge each system contributes to Fan \& Lan's comparison and what remains absent.

\paragraph{13.4.1 Runtime boundaries that could host recovery}\label{runtime-boundaries-that-could-host-recovery}

\begin{itemize}
\tightlist
\item
  LangGraph: state graph, checkpoint, interrupt and resume boundary.
\item
  OpenHands: execution/workflow state and agent-execution boundary.
\item
  Magentic-One: orchestration boundary relevant to multi-agent task execution.
\item
  Shepherd: reversible first-class execution traces, including inspect, transform, fork and replay.
\item
  These are candidate insertion boundaries, not proof of B2. A checkpoint by itself is persistence at most; it becomes a Soar-related D3 control correspondence only when a typed failure trigger opens a bounded child state and controlled return semantics are present.
\end{itemize}

\paragraph{13.4.2 Diagnosis and recovery baselines}\label{diagnosis-and-recovery-baselines}

\begin{itemize}
\tightlist
\item
  \textbf{PALADIN:} Fan \& Lan describe it as classifying tool errors and retrieving recovery actions. The preserved edge is that failure type changes which recovery action is selected.
\item
  \textbf{AgentDebugX:} Fan \& Lan describe a Detect--Attribute--Recover--Rerun loop. The preserved edge links failure observability and attribution to repair and rerun.
\item
  Explain that these are diagnostic/recovery baselines and should not be presented as implementing the protected substate, typed return and safe compilation chain unless the source says so.
\end{itemize}

\paragraph{13.4.3 Reversible execution and cross-episode learning baselines}\label{reversible-execution-and-cross-episode-learning-baselines}

\begin{itemize}
\tightlist
\item
  \textbf{Shepherd:} reversible, forkable and replayable first-class execution traces; useful as the closest protected-branch baseline.
\item
  \textbf{AgentHER and Reflexion:} failure experience can affect later behavior; useful as learning baselines, but not equivalent to compiling the dependency structure of a scoped Soar substate.
\item
  Distinguish in-run recovery and branching from cross-episode feedback.
\end{itemize}

\paragraph{13.4.4 A source-faithful comparison table}\label{a-source-faithful-comparison-table}

Planned table columns: system/case; state object; trigger; preserved edge; what is not established; Fan \& Lan's provisional E/D record where the paper gives one.

\begin{itemize}
\tightlist
\item
  PALADIN and AgentDebugX: functional recovery/diagnosis baselines; Fan \& Lan's Appendix C-B reports the typed-recovery records as provisional E3/D3.
\item
  Shepherd: protected-branch record; Appendix C-C reports E2/D3.
\item
  AgentHER and Reflexion: failure-learning records; Appendix C-D reports E3/D4 for the narrow feedback edge.
\item
  Do not assign one E/D code to the entire Soar-to-modern mapping or to B2 as a bundle.
\item
  Refer back to Chapter 11 for the definitions of E1--E4 and D0--D4. Only explain what is needed to interpret these particular record labels.
\end{itemize}

\subsubsection{13.5 B2 as a composed runtime hypothesis}\label{b2-as-a-composed-runtime-hypothesis}

This section reconstructs the five-step proposal in Fan \& Lan and then makes the dependencies visible. It must not pretend the proposal is already a validated architecture.

\paragraph{13.5.1 Step 1 --- classify the failure}\label{step-1-classify-the-failure}

\begin{itemize}
\tightlist
\item
  An execution event crosses a monitored boundary: a tool call, agent action, validation step or workflow transition fails or cannot establish a required postcondition.
\item
  The execution ledger records the event and classifies the failure type; root-cause attribution is recorded separately from the raw symptom.
\item
  Classification includes evidence and uncertainty. Where the type cannot be established, the safe result may be abstention or escalation rather than speculative recovery.
\end{itemize}

\paragraph{13.5.2 Step 2 --- preserve the parent state}\label{step-2-preserve-the-parent-state}

\begin{itemize}
\tightlist
\item
  Record the parent task/state identifier, version or hash, goal, relevant invariants, completed verified work, pending action and dependencies.
\item
  Freeze or otherwise protect the relevant parent view while recovery proceeds.
\item
  Define which observations may be added, which state changes are local to the child, and which external side effects cannot be rolled back by restoring an internal snapshot.
\item
  Explicitly label the concrete schema and snapshot policy as a book-level engineering proposal.
\end{itemize}

\paragraph{13.5.3 Step 3 --- open a bounded recovery subgraph}\label{step-3-open-a-bounded-recovery-subgraph}

\begin{itemize}
\tightlist
\item
  Create a child recovery scope carrying the typed failure, parent reference, evidence and allowed recovery operations.
\item
  Let the child have its own local decision process, but do not let it silently rewrite the parent goal or treat unverified proposals as parent results.
\item
  Bound attempts and permitted operations. Mention deadline/token/tool limits only as local guardrails; defer cross-stage resource allocation and interruption policy to Chapter 16 (B6).
\item
  Record nested recovery only when a new, separately typed impasse warrants it; avoid unbounded recursive retries.
\end{itemize}

\paragraph{13.5.4 Step 4 --- return and validate a typed resolution}\label{step-4-return-and-validate-a-typed-resolution}

\begin{itemize}
\tightlist
\item
  Return a typed outcome, not an unstructured assistant paragraph.
\item
  The resolution should carry: outcome type; evidence; affected resources; parent version; dependency references; completed/failed operations; postconditions; and whether the result is safe to apply.
\item
  The parent checks that it is still the same relevant state, that dependencies remain valid, and that the result satisfies the expected contract.
\item
  If the parent has changed or the evidence is stale, the result is rejected, re-evaluated or the branch regenerated; it is not silently merged.
\end{itemize}

\paragraph{13.5.5 Step 5 --- compile a reusable recovery rule}\label{step-5-compile-a-reusable-recovery-rule}

\begin{itemize}
\tightlist
\item
  Extract only the conditions that made the resolution valid and the validated recovery action/result.
\item
  Express those conditions as explicit guards or typed preconditions, not merely a semantic similarity to a past failure.
\item
  Store provenance, confidence/validation status, source versions, tool/environment identity, rule version and invalidation conditions.
\item
  Run candidate rules in a controlled validation path before promotion; demote or disable them after counterexamples or dependency changes.
\item
  Clarify that this is a proposed modern analogue of dependency-based chunking, not a claim that Fan \& Lan supply a complete compiler algorithm or that all failed runs should be compiled.
\end{itemize}

\paragraph{13.5.6 The complete lifecycle and its authority boundaries}\label{the-complete-lifecycle-and-its-authority-boundaries}

\begin{itemize}
\tightlist
\item
  Show the chain from parent operation to typed detection, parent preservation, child recovery, typed return, parent validation and optional compilation.
\item
  Identify the authoritative controller at every transition. Model-generated hypotheses may propose a diagnosis or recovery action; the runtime governor decides whether the branch opens, whether a result is accepted and whether a rule is promoted.
\end{itemize}

\subsubsection{13.6 Product X worked example: a payment retry with an unknown outcome}\label{product-x-worked-example-a-payment-retry-with-an-unknown-outcome}

\textbf{{[}Book example{]}} Reuse Product X from Chapter 12 to maintain continuity, while marking every added scenario detail as illustrative.

\begin{itemize}
\tightlist
\item
  Parent goal: resolve the payment attempt without creating a duplicate charge; preserve the customer request and verified transaction facts.
\item
  Event: the payment retry returns a timeout. A timeout alone does not prove that the provider rejected the transaction.
\item
  Typed classification: \texttt{payment\_outcome\_unknown}, supported by the timeout event and the absence of a confirmed status response.
\item
  Parent preservation: record the parent task version and invariant that no second charge may be initiated until the first attempt's status is known.
\item
  Child recovery: perform only allowed status checks using the original idempotency/transaction reference; do not initiate a fresh payment in the recovery branch.
\item
  Typed resolution alternatives:

  \begin{itemize}
  \tightlist
  \item
    \texttt{settled}: return verified provider evidence; parent records completion and does not retry.
  \item
    \texttt{not\_settled}: return verified evidence; parent may consider the permitted retry path.
  \item
    \texttt{still\_unknown}: return uncertainty; parent escalates or waits rather than guessing.
  \end{itemize}
\item
  Compilation candidate: after repeated validated episodes, a narrow rule may be proposed: when the same provider/API version reports the same timeout class and a status query proves the original attempt is not settled, use the verified retry procedure with the same idempotency reference.
\item
  Guard against unsafe overgeneralization: a timeout with unknown settlement is not equivalent to a confirmed failed transaction; provider version, idempotency semantics and postconditions matter.
\item
  Include a counterexample where the provider's status result changes while the branch is open. The child result must be invalidated/rechecked rather than returned as though its parent context were unchanged.
\item
  Make the reader identify which parts are source mechanism, which are Fan \& Lan's B2 proposal, and which are the book's engineering example.
\end{itemize}

\subsubsection{13.7 Apply the seven-field mechanism representation to B2}\label{apply-the-seven-field-mechanism-representation-to-b2}

Use the exact tuple from Chapter 10, with no new fields and no redefinition. The purpose is to summarize the proposal, not to repeat the entire Chapter 10 lesson.

\begin{longtable}[]{@{}
  >{\raggedright\arraybackslash}p{(\columnwidth - 4\tabcolsep) * \real{0.3333}}
  >{\raggedright\arraybackslash}p{(\columnwidth - 4\tabcolsep) * \real{0.3333}}
  >{\raggedright\arraybackslash}p{(\columnwidth - 4\tabcolsep) * \real{0.3333}}@{}}
\toprule\noalign{}
\begin{minipage}[b]{\linewidth}\raggedright
Field
\end{minipage} & \begin{minipage}[b]{\linewidth}\raggedright
Soar historical mechanism
\end{minipage} & \begin{minipage}[b]{\linewidth}\raggedright
B2 runtime hypothesis / proposed design
\end{minipage} \\
\midrule\noalign{}
\endhead
\bottomrule\noalign{}
\endlastfoot
State substrate (S) & Working-memory states, operators, preferences, substates/superstates, results and dependency information. & Execution ledger, parent-state version, typed failure, recovery subgraph trace, typed result and compiled recovery rule. \\
Control locus (C) & Soar's decision procedure and production/operator-selection machinery. & A runtime recovery governor with authority to classify, fork, accept/reject results and promote/reject compilation candidates. \\
Transition trigger (T) & Operator-selection impasse; successful resolution or dependency/context change. & Typed execution failure or attributed cause; child resolution; parent validation; later matching invocation. \\
Persistence boundary (P) & Child-local state is removed on resolution; linked results survive; learned chunks persist as productions. & Child branch is scoped to the recovery episode; parent state remains authoritative; validated rules persist with versions and provenance. \\
Failure semantics (F) & Typed impasse, resolution, elimination and regeneration. & Misclassification, unresolved recovery, stale parent/dependency, rejected result or unsafe compiled rule. \\
Learning operator (L) & Chunking compiles dependency conditions and result into a production after successful impasse resolution. & Proposed condition-safe compilation and later outcome-based rule validation/update. Keep AgentHER/Reflexion's separate learning edge distinct. \\
Resource/governance policy (R) & Only report what the source documents; do not fill unknown architectural resource details. & Bound recovery attempts/operations locally; cross-stage budget, risk and interruption authority belong to B6. \\
\end{longtable}

\subsubsection{13.8 Failure modes, safety and authority}\label{failure-modes-safety-and-authority}

\paragraph{13.8.1 Diagnosis errors}\label{diagnosis-errors}

\begin{itemize}
\tightlist
\item
  False diagnosis opens the wrong recovery path; missed diagnosis falls back to generic retries; ambiguous cause should not be promoted to a durable type without evidence.
\item
  Measure per-type precision/recall and the cost of wrong recovery selection.
\end{itemize}

\paragraph{13.8.2 Parent-state corruption and stale context}\label{parent-state-corruption-and-stale-context}

\begin{itemize}
\tightlist
\item
  A child mutates parent data, merges speculative state, or returns a result calculated against an outdated parent version.
\item
  Use parent-version checks and dependency-aware invalidation as the proposed control response.
\item
  Explain that an internal snapshot cannot reverse an external side effect already performed by a tool.
\end{itemize}

\paragraph{13.8.3 Unsafe compilation and overgeneralization}\label{unsafe-compilation-and-overgeneralization}

\begin{itemize}
\tightlist
\item
  A recovery rule is learned from a one-off incident, an incorrect attribution, missing dependencies or a context-specific result.
\item
  The learned guard is too broad, the tool contract changes, or the rule remains active after a dependency changes.
\item
  Require provenance, narrow applicability, validation, expiry/invalidation and rollback as engineering safeguards.
\end{itemize}

\paragraph{13.8.4 Endless recovery, retry storms and authority confusion}\label{endless-recovery-retry-storms-and-authority-confusion}

\begin{itemize}
\tightlist
\item
  A recovery branch recursively opens new branches without a termination condition.
\item
  The model both proposes and authorizes a state-changing recovery.
\item
  A branch crosses the permitted scope or repeats an external side effect.
\item
  Keep these concerns distinct from B6's full resource/uncertainty/interruption control policy, but show why B2 still needs a local maximum depth and escalation path.
\end{itemize}

\subsubsection{13.9 Evaluation: test the full coupling, not only whether a retry worked}\label{evaluation-test-the-full-coupling-not-only-whether-a-retry-worked}

\paragraph{13.9.1 Required trace and event record}\label{required-trace-and-event-record}

Use the fields specified by Fan \& Lan's Appendix D:

\begin{itemize}
\tightlist
\item
  failure type;
\item
  parent hash or stable parent-state reference;
\item
  substate/recovery-branch trace;
\item
  returned resolution;
\item
  compiled rule and its applicability record.
\end{itemize}

Add proposed diagnostic fields only when clearly labelled as engineering additions: classifier evidence, parent version, dependency set, branch budget, side-effect ledger, result validation and rule version.

\paragraph{13.9.2 Baselines and controlled comparisons}\label{baselines-and-controlled-comparisons}

\begin{itemize}
\tightlist
\item
  Simple retry / generic retry loop.
\item
  Typed diagnosis with recovery selection but no protected child state.
\item
  Reversible branch/replay without typed impasse and typed return.
\item
  Cross-episode reflection or hindsight learning without dependency-safe compilation.
\item
  Full B2 composition.
\item
  A human-supervised or conservative escalation path where the task's risk makes automatic recovery inappropriate.
\end{itemize}

Keep each comparison tied to the specific edge it isolates; do not imply every baseline is directly implemented by a single named product.

\paragraph{13.9.3 Primary, safety and efficiency measures}\label{primary-safety-and-efficiency-measures}

\begin{itemize}
\tightlist
\item
  Recovery success by failure type.
\item
  Repeated-failure rate and transfer to a later similar case.
\item
  Parent-state corruption/invariant-violation rate.
\item
  Unsafe rule-reuse rate and harmful compilation rate.
\item
  Diagnosis/recovery latency.
\item
  Recovery operations, tool calls and cost.
\item
  Rejection/escalation precision for unresolved or ambiguous cases.
\end{itemize}

Fan \& Lan's explicit Table X measures include recovery by type, repeated-failure rate and transfer as primary outcomes; diagnosis latency, state corruption and unsafe reuse as safety/efficiency outcomes. The chapter should preserve those labels and may add operational definitions as book-level measurement design.

\paragraph{13.9.4 Required ablations and rejection conditions}\label{required-ablations-and-rejection-conditions}

\begin{itemize}
\tightlist
\item
  Remove the new state object while retaining the rest of the controller.
\item
  Remove the authoritative control edge but retain logging, to test whether the runtime transition rather than instrumentation produces the effect.
\item
  Remove the learning update while retaining the immediate runtime recovery transition, to separate D3 control value from D4 adaptation.
\item
  Reject the B2 claim if the retry baseline is non-inferior under the preregistered margin, if parent-state corruption increases, if compilation causes harmful reuse, or if repeated failures/transfer do not improve.
\item
  Follow Fan \& Lan's Appendix D discipline: hold base model, prompt, tools, environment and task budget fixed; preregister task population, experimental unit, versions, stochastic replications, sample size, metrics, margins/tests, exclusions and stopping rule.
\end{itemize}

\subsubsection{13.10 Boundaries with other bundles and handoff}\label{boundaries-with-other-bundles-and-handoff}

\begin{itemize}
\tightlist
\item
  \textbf{B1 (Chapter 12):} memory-item accessibility and context selection; it may supply prior information to recovery but does not define the failure transition.
\item
  \textbf{B2 (this chapter):} typed failure, protected substate, typed return and resolution compilation inside an execution episode.
\item
  \textbf{B3 (Chapter 14):} competition for a bounded workspace and broadcast/admission learning; it may receive recovery outputs as candidate content but is not the recovery protocol.
\item
  \textbf{B4 (Chapter 15):} intention and reconsideration authority; it supplies the goal boundary that recovery must not silently rewrite.
\item
  \textbf{B5 (Chapter 17):} general skill governance. Fan \& Lan keep B2 separate because B2 is an in-episode recovery transition while B5 concerns general skill invocation across episodes/tasks; the compilation of a narrow recovery rule does not turn B2 into the entire skill-governance problem.
\item
  \textbf{B6 (Chapter 16):} uncertainty, resources, interruption and stopping. B6 governs wider risk/budget authority; B2 retains local recovery bounds and escalation.
\item
  \textbf{Chapter 18:} integrates the dependencies after the five residual bundles and closed B5 case have been examined individually.
\end{itemize}

\subsubsection{Closing --- Recap and review questions}\label{closing-recap-and-review-questions}

The recap should return to the core sequence: typed impasse/failure → scoped child recovery → validated result return → dependency-safe compilation. It should emphasize the boundary between the native Soar semantics, Fan \& Lan's B2 proposal, and the book's proposed engineering contract.

Suggested review questions:

\begin{enumerate}
\def\labelenumi{\arabic{enumi}.}
\tightlist
\item
  Why is a Soar impasse not the same thing as a tool exception?
\item
  What distinguishes a tie impasse from a conflict or constraint-failure impasse?
\item
  What does the \texttt{superstate} relationship contribute to a substate?
\item
  What is a result in Soar, and why is it different from arbitrary child output?
\item
  Explain the difference between resolving, eliminating and regenerating an impasse.
\item
  What does the Goal Dependency Set track, and why can a dependency change invalidate child reasoning?
\item
  Which conditions and action are used in a chunk, and why are they variablized?
\item
  Why does a checkpoint alone not establish a typed substate?
\item
  What do PALADIN, AgentDebugX, Shepherd, AgentHER and Reflexion contribute to the B2 comparison, and what does each leave unestablished?
\item
  Why must a modern recovery branch return a typed result tied to the parent version?
\item
  Give an example of an unsafe generalization from a successful recovery episode.
\item
  Which B2 metrics test immediate recovery, which test state safety, and which test learning/transfer?
\item
  Explain why B2 and B5 remain separate even though B2 may compile a reusable recovery rule.
\item
  Which parts of B2 depend on the intention authority in B4 or the resource/uncertainty controls in B6?
\end{enumerate}

\subsection{4. Planned figures and diagrams}\label{planned-figures-and-diagrams}

All conceptual diagrams are to be authored as Mermaid \texttt{.mmd} files under \texttt{Tex/diagrams/source/part2/chapter13/}, rendered by the shared pinned Mermaid CLI in CI, and placed in LaTeX only through \texttt{\textbackslash{}BookDiagram}. Every \texttt{.mmd} file must include \texttt{accTitle} and \texttt{accDescr}, use the shared theme and have a matching caption and nearby explanatory prose. No TikZ, ASCII diagrams, hand-positioned LaTeX nodes or chapter-specific drawing packages.

\begin{enumerate}
\def\labelenumi{\arabic{enumi}.}
\item
  \textbf{Figure 13.1 --- Soar operator decision and typed impasse cycle.} Show the parent state, operator proposals/preferences, selection and the normal operator branch versus a typed impasse branch; show child substate processing, a result linked to the parent, resolution and optional chunking. Include separate paths for elimination/regeneration so the diagram does not imply every disappearing impasse is learned.

  Source file: \path{soar-impasse-substate-cycle.mmd}.
\item
  \textbf{Figure 13.2 --- The four B2 invariants and their dependency chain.} Show typed diagnosis establishing recovery scope; protected parent state and bounded child substate; typed result return and parent validation; dependency-aware compilation into a reusable recovery rule. Mark this as Fan \& Lan's proposed bundle, not an already implemented standard.

  Source file: \path{b2-four-invariants.mmd}.
\item
  \textbf{Figure 13.3 --- Current evidence covers separate B2 edges.} Map PALADIN/AgentDebugX to diagnosis and recovery selection, Shepherd to reversible execution branches, and AgentHER/Reflexion to later behavior change. Show that the integrated typed impasse--protected substate--typed return--safe compilation chain remains the residual.

  Source file: \path{b2-modern-baseline-map.mmd}.
\item
  \textbf{Figure 13.4 --- Product X recovery contract.} Trace a timeout through an unknown-payment classification, preserved parent, read-only child status check, typed status alternatives, parent validation and a narrowly guarded compilation candidate. Caption and prose must mark this as a book-created illustration.

  Source file: \path{productx-b2-recovery-example.mmd}.
\end{enumerate}

\subsection{5. Shared glossary and index plan}\label{shared-glossary-and-index-plan}

Do not create duplicate glossary entries. Before drafting, check the shared glossary and add only missing entries. Candidate entries for Chapter 13:

\begin{itemize}
\tightlist
\item
  impasse;
\item
  tie impasse;
\item
  conflict impasse;
\item
  constraint-failure impasse;
\item
  no-change impasse;
\item
  operator;
\item
  preference;
\item
  substate;
\item
  superstate;
\item
  subgoal result;
\item
  Goal Dependency Set (GDS);
\item
  chunking;
\item
  parent-state snapshot;
\item
  typed resolution;
\item
  recovery subgraph;
\item
  dependency condition.
\end{itemize}

Use \texttt{\textbackslash{}gls} on first meaningful use in the chapter and \texttt{\textbackslash{}index\{...\}} consistently for Soar, operator selection, impasses, substates, superstates, subgoal results, GDS, chunking, recovery branches, typed results, compilation and transfer. Definitions should explain the term at the level needed by a beginner and must not imply that modern design terms are native Soar terminology.

\subsection{6. Teaching boxes and running example plan}\label{teaching-boxes-and-running-example-plan}

Use the shared semantic boxes selectively:

\begin{itemize}
\tightlist
\item
  \texttt{learningbox}: learning outcomes.
\item
  \texttt{takeawaybox}: B2 thesis / the full typed-failure-to-safe-compilation chain.
\item
  \texttt{notebox}: the distinction between Soar's superstate dependency model and a modern immutable parent snapshot.
\item
  \texttt{warningbox}: a tool timeout is not proof that an external side effect did not occur; restoring an internal checkpoint does not undo external effects.
\item
  \texttt{tipbox}: how to read the modern-baseline comparison as separate preserved edges, not as system rankings.
\item
  \texttt{chapterrecap}: concise mechanism recap.
\item
  \texttt{reviewquestions}: exercises and conceptual questions.
\end{itemize}

Keep ordinary prose out of teaching boxes. Use the Product X payment incident for continuity with Chapter 12 and mark all invented incident details \texttt{{[}Book\ example{]}}; mark proposed schemas and controller rules \texttt{{[}Engineering\ decision{]}}.

\subsection{7. Source and attribution plan}\label{source-and-attribution-plan}

\subsubsection{Primary theoretical source --- Fan \& Lan (2026)}\label{primary-theoretical-source-fan-lan-2026}

\begin{itemize}
\tightlist
\item
  \textbf{B2 hypothesis and exact proposed chain:} §VI-C, B2; quote/paraphrase the execution ledger, parent freeze, bounded recovery subgraph, typed resolution and reusable recovery rule.
\item
  \textbf{Why the bundle consists of these four residual invariants:} §VI-B, Table V, B2 row: typed diagnosis, protected substate, typed return and resolution compilation; explain why B2 remains separate from B5.
\item
  \textbf{Historical Soar mechanism:} §III-A, Soar; Figure 2(b); historical architecture ledger/Table VIII and Appendix A where applicable.
\item
  \textbf{Historical-to-modern mapping:} §V-A, ``Soar to typed recovery and experience learning''; §V-B, ``Failure changes from a control state to recovery content''; Table IV.
\item
  \textbf{Modern baseline descriptions:} §IV-B--IV-C; §V-A Soar mapping; Appendix C-B (R2a), C-C (R2b), C-D (R3); Appendix F-B/Table XIII.
\item
  \textbf{E/D status for the particular records:} Appendix C-B reports typed recovery as provisional E3/D3; Appendix C-C reports Shepherd's protected-branch mapping as E2/D3; Appendix C-D reports the AgentHER/Reflexion learning edge as E3/D4. Refer to Chapter 11 for full definitions and do not create a B2-wide combined code.
\item
  \textbf{Checkpoint boundary:} Appendix E-B, ``Checkpoint versus typed impasse''; a checkpoint is persistence at D2, while a Soar-related D3 claim needs a typed failure trigger, bounded child state and controlled return semantics.
\item
  \textbf{Experiment protocol and falsification:} Appendix D, Table X, B2 row and general preregistration/ablation requirements; §VI-C, B2 failure conditions; §VI-C Table VI, B2 row.
\item
  \textbf{Cross-bundle dependencies:} §VI-D and Figure 8; B4 scopes recovery, B6 governs risk/resources, and B2 can produce verified resolutions for a GraSP-like governor.
\item
  \textbf{Validity boundary:} §VII-C and Appendix F-C. The paper reports no completed experiment comparing the residual bundles against their baselines; present B2 as a hypothesis/research agenda, not a demonstrated performance result.
\end{itemize}

\subsubsection{Historical mechanism cross-check}\label{historical-mechanism-cross-check}

The Soar Architecture manual cited by Fan \& Lan as reference {[}9{]} is used to verify the named impasse types, creation of substates, result semantics, resolution/elimination/regeneration, Goal Dependency Set and chunking. Cite these source-specific details with \autocite{soar2025architecture}, alongside \FanLanCite{...} for Fan \& Lan's reconstruction and B2 analysis. This is not a third theoretical foundation; it is the historical architecture source already cited within the primary paper.

\subsubsection{Citation discipline}\label{citation-discipline}

Use the shared \texttt{\textbackslash{}FanLanCite\{...\}} macro for the review's framing, mapping, hypothesis and validity claims. Use \texttt{\textbackslash{}autocite\{soar2025architecture\}} only for the Soar-specific facts checked in that manual. Other named modern systems may appear only as cases that Fan \& Lan explicitly include; do not turn them into separate theory chapters or assert mechanisms beyond the review's stated descriptions. Keep claims from the book-created Product X scenario and proposed runtime contracts labelled as such.

\subsection{8. Explicit exclusions and chapter boundaries}\label{explicit-exclusions-and-chapter-boundaries}

\begin{itemize}
\tightlist
\item
  Do not treat every failed tool call, timeout or validation error as a Soar impasse.
\item
  Do not suggest that Soar substates are simply OS processes, agent forks or immutable snapshots.
\item
  Do not claim that the superstate is literally frozen in the same implementation sense as Fan \& Lan's proposed parent-state freeze.
\item
  Do not imply that a child result is any arbitrary text returned from a model; explain Soar's result linkage and the proposed typed-resolution contract separately.
\item
  Do not say that an impasse always resolves. Preserve elimination and regeneration as distinct outcomes.
\item
  Do not claim chunking occurs after every failed attempt or after impasse elimination/regeneration; connect it to successful resolution.
\item
  Do not equate Soar chunking with reflection memories, trajectory replay, retrieval-augmented prompting or generic skill storage.
\item
  Do not claim PALADIN, AgentDebugX, Shepherd, AgentHER or Reflexion independently implements the full B2 chain.
\item
  Do not treat E/D codes as one score or assign a single code to the whole B2 bundle; Chapter 11 owns their definitions.
\item
  Do not turn the book's proposed resolution fields, controller authority, compile guards or Product X incident into statements attributed to Soar or Fan \& Lan.
\item
  Do not teach B3 workspace control, B4 intention governance, B5 skill governance or B6 global uncertainty/resource control in depth; point to Chapters 14--17.
\item
  Do not claim B2 has been experimentally validated. The primary paper describes a research hypothesis and a proposed falsification protocol, not a completed comparative intervention.
\end{itemize}

\subsection{9. Proposed depth and acceptance criteria}\label{proposed-depth-and-acceptance-criteria}

\textbf{Expected finished-chapter depth:} approximately 30--40 pages, with no fixed page-count target. The chapter should be long enough to explain the Soar mechanism accurately, reconstruct the B2 hypothesis, compare partial modern baselines, develop one worked example and teach the evaluation protocol without compressing the distinctions.

The chapter draft will be ready for review only when it:

\begin{enumerate}
\def\labelenumi{\arabic{enumi}.}
\tightlist
\item
  Starts from the problem of safe recovery and introduces terms only when they are needed.
\item
  Teaches Soar's operator/preferences decision point before explaining typed impasses.
\item
  Describes the four impasse types and does not conflate them with generic execution exceptions.
\item
  Explains the substate/superstate relationship, result semantics and the distinction between resolution, elimination and regeneration.
\item
  Explains GDS dependency tracking and the dependency basis of chunking accurately.
\item
  Separates native Soar behavior from Fan \& Lan's proposed modern execution-ledger design.
\item
  Reconstructs all four B2 invariants and their dependencies without changing the hypothesis into an asserted implementation standard.
\item
  Compares PALADIN, AgentDebugX, Shepherd, AgentHER and Reflexion only at the level supported by Fan \& Lan.
\item
  Treats LangGraph, OpenHands, Magentic-One and Shepherd as insertion/state/branching boundaries rather than proof of complete B2.
\item
  Uses source-specific E/D statuses without redefining or collapsing the axes from Chapter 11.
\item
  Provides an internally coherent Product X example with an explicit \texttt{{[}Book\ example{]}} label.
\item
  Applies the seven-field tuple exactly as introduced in Chapter 10.
\item
  Preserves Fan \& Lan's B2 trace, metrics, ablations and falsification conditions.
\item
  Distinguishes B2 from the in-scope/out-of-scope neighboring bundles, especially B5 and B6.
\item
  Uses the shared preamble, macros, semantic box policy, footnote citations, glossary/index, Mermaid sources and chapter-local artifact conventions.
\item
  Is built with XeLaTeX/Biber, passes scaffold/diagram/PDF checks and has its rendered pages and actual CI-generated diagrams visually inspected before final chapter PDF release.
\end{enumerate}

\subsection{10. Approval checkpoint}\label{approval-checkpoint}

\textbf{Status: Awaiting explicit author approval.} This document is the proposed Chapter 13 outline only. No Chapter 13 prose has been drafted. Approval of Chapter 12 or the batch plan does not automatically approve this chapter outline. After approval, the next stage is the Chapter 13 prose draft, followed by source-fidelity review, Mermaid rendering, XeLaTeX/Biber build, PDF inspection and repository integration.

\begin{center}\rule{0.5\linewidth}{0.5pt}\end{center}

\subsection{Primary reference}\label{primary-reference}

Fan, H., \& Lan, Z. (2026). \emph{From Cognitive Architectures to Language Agents: A Mechanism-Level Review of Lineage, Convergence, and Migration Gaps}. arXiv:2607.23942v1, submitted 27 July 2026. Principal anchors: §III-A (Soar); §V-A (Soar mapping); §V-B (failure translation); §VI-B, Table V; §VI-C, B2 and Table VI; §VI-D; Figure 8; Appendices C-B, C-C, C-D, D, E-B and F-B/F-C.

\FanLanCite{\S III-A, Soar; \S V-A--V-B; \S VI-B, Table V; \S VI-C, B2 and Table VI; \S VI-D, Figure 8; Appendices C-B--C-D, D, E-B and F-B--F-C}

\cleardoublepage
\printbibliography[heading=bibintoc,title={Historical mechanism cross-check}]
\end{document}
